% =====================================================================
% comum/theoremstyles.tex — ambientes de teorema (amsthm) + cleveref
% =====================================================================
% Requer: amsthm (pacotes/matematica.tex) e cleveref (pacotes/referencias.tex).
%
% Ambientes disponíveis:
%   estilo plain      : theorem, proposition, lemma, corollary
%   estilo definition : definition, axiom, example, examples, observation
%   sem número        : theorem*, definition*, example*, observation*
%
% Configuração (defina ANTES de carregar este arquivo):
%   \newcommand{\comumNumeracao}{chapter}  % Teorema 2.1  (livro/TCC)
%   \newcommand{\comumNumeracao}{section}  % Teorema 2.1  (artigo)
%   (não definir)                           % Teorema 1, 2, ...
%
%   \newcommand{\comumContadorUnico}{}     % todos os ambientes dividem um
%                                          % contador (Def. 2.1, Teo. 2.2, ...)
%                                          % padrão: um contador por ambiente
%
% Uso com cleveref: \cref{teo:x} -> "teorema 2.1"; \Cref{teo:x} -> "Teorema 2.1".
% Prefixos de rótulo sugeridos: teo:, prop:, lem:, cor:, def:, ax:, ex:, obs:
% =====================================================================

\providecommand{\comumNumeracao}{}
\ifdefined\comumContadorUnico
  \RequirePackage{aliascnt}
\fi

\makeatletter
% \comum@novoteorema{ambiente}{Nome}
\newcommand{\comum@novoteorema}[2]{%
  \ifdefined\comumContadorUnico
    % o primeiro ambiente (theorem) cria o contador; os demais o reutilizam
    \@ifundefined{c@theorem}%
      {\comum@teoremasimples{#1}{#2}}%
      {% aliascnt: contador "apelido" de theorem, para o cleveref saber
       % que o rótulo é uma definição/lema/... e não um teorema
       \newaliascnt{#1}{theorem}%
       \newtheorem{#1}[#1]{#2}%
       \aliascntresetthe{#1}}%
  \else
    \comum@teoremasimples{#1}{#2}%
  \fi
}
\newcommand{\comum@teoremasimples}[2]{%
  \if\relax\detokenize\expandafter{\comumNumeracao}\relax
    \newtheorem{#1}{#2}%
  \else
    \newtheorem{#1}{#2}[\comumNumeracao]%
  \fi
}

% ---------- Estilo plain (itálico): resultados ----------
\theoremstyle{plain}
\comum@novoteorema{theorem}{Teorema}
\comum@novoteorema{proposition}{Proposição}
\comum@novoteorema{lemma}{Lema}
\comum@novoteorema{corollary}{Corolário}
\newtheorem*{theorem*}{Teorema}

% ---------- Estilo definition (reto): definições e exemplos ----------
\theoremstyle{definition}
\comum@novoteorema{definition}{Definição}
\comum@novoteorema{axiom}{Axioma}
\comum@novoteorema{example}{Exemplo}
% examples compartilha o contador de example
\newtheorem{examples}[example]{Exemplos}
\comum@novoteorema{observation}{Observação}
\newtheorem*{definition*}{Definição}
\newtheorem*{example*}{Exemplo}
\newtheorem*{observation*}{Observação}

% ---------- Fim de demonstração ----------
\renewcommand{\qedsymbol}{$\blacksquare$}

% ---------- Nomes para \cref (minúsculo) e \Cref (maiúsculo) ----------
\crefname{theorem}{teorema}{teoremas}
\Crefname{theorem}{Teorema}{Teoremas}
\crefname{proposition}{proposição}{proposições}
\Crefname{proposition}{Proposição}{Proposições}
\crefname{lemma}{lema}{lemas}
\Crefname{lemma}{Lema}{Lemas}
\crefname{corollary}{corolário}{corolários}
\Crefname{corollary}{Corolário}{Corolários}
\crefname{definition}{definição}{definições}
\Crefname{definition}{Definição}{Definições}
\crefname{axiom}{axioma}{axiomas}
\Crefname{axiom}{Axioma}{Axiomas}
\crefname{example}{exemplo}{exemplos}
\Crefname{example}{Exemplo}{Exemplos}
\crefname{examples}{exemplos}{exemplos}
\Crefname{examples}{Exemplos}{Exemplos}
\crefname{observation}{observação}{observações}
\Crefname{observation}{Observação}{Observações}
\makeatother
